\chapter{Conclusions} \label{Chapter:conclusions}

We have presented a way to correct classical binding energies and the steps leading up to the final method.  These steps have shown several important problems with quantum corrections.  Chapter XXX shows the issue of using a SSP approach to correcting free energies, namely that large ``jumps'' are present within the free energy.  These ``jumps' are as a result of energy outliers at the low energy tail of the energy difference distribution.  This distribution has a large range, even for the simple systems used within this Thesis.  If further examination is performed on the outliers within the energy difference distribution, it is found that they are not outliers in either the classical or quantum energy distributions.  However, they are on different sides of the distributions.  This highlights the inexact overlap of the configurational space.

The issue of inexact overlap was then further investigated by perturbing between two classical potentials with different charges.  It was found that by adjusting the overlap slightly the energy does not converge not only between runs, but the forward and reverse process too.  This showed the sensitivity of the Zwanzig equation and led to the examination of post-processing methods with an aim to ``smooth'' over the low energy tail.  First a Gaussian distribution was fitted to the energy difference.  Numerical integration showed much improved convergence on the direct SSP approach, but still yielded a difference of over 6 kcal/mol between runs.  Additionally an analytical approach was applied, which emphasized the sensitivity of the exponents on the convergence of the free energy.  Specifically, the variance or $\alpha$ exponent has the largest effect on the free energy.  Again this is due to the low energy tail of the distribution, by only allowing data from ten thousandth height of the Gaussian the free energy converges.  However, there is no scientific reason to ignore part of the data, which shows that no post-processing method can be applied to increase the convergence.  This led to the importance of finding an accurate representation of the QM configurational space.






The next steps for this project would be to apply the final method to a more real test system.  This should be much closer to a protein-ligand binding site, but much smaller, for example, cyclodextrin bound to a variety of small organic ligands.  This would provide a non-uniform polarity around the ligand and should show the real benefit this method can have on complex systems.